%% kapa: exam=1 round=29 year=2018 period=2 form=A subject=accounting subject_ko=회계학 seq=33 number=73
%% source: Q-Net 감정평가사 기출문제 (cst00309) artlSeq=5202808 file=extracted/1008/5202808/2018년 제29회 감정평가사 1차시험 기출문제/2018년 제29회 감정평가사 1차시험 기출문제/2교시 A형.pdf
%% conversion: kordoc 4.24.51 -> tools/gen1.py (auto); review=proofread-sonnet; flags=none
\begin{question}{73}
실제개별원가계산제도를 사용하는 (주)감평의 20x1년도 연간 실제 원가는 다음과 같다.
\begin{framedtext}
\begin{tabular}{@{}p{0.2\linewidth}p{0.2\linewidth}p{0.2\linewidth}l@{}}
직접재료원가 & ￦4,000,000 & 직접노무원가 & ￦5,000,000 \\
제조간접원가 & ￦1,000,000 & & \\
\end{tabular}
\end{framedtext}
(주)감평은 20x1년 중 작업지시서 \#901을 수행하였는데 이 작업에 320시간의 직접노무시간이 투입되었다. (주)감평은 제조간접원가를 직접노무시간을 기준으로 실제배부율을 사용하여 각 작업에 배부한다. 20x1년도 실제 총직접노무시간은 2,500시간이다. (주)감평이 작업지시서 \#901에 배부하여야 할 제조간접원가는? \par
\begin{choices}
  \item ￦98,000
  \item ￦109,000
  \item ￦128,000
  \item ￦160,000
  \item ￦175,000
\end{choices}
\end{question}
\answer{3}
